\documentclass[10pt,a4paper]{amsart}
\usepackage[margin=19mm]{geometry}
\usepackage{amsmath,amssymb,mathtools}
\usepackage[scr=rsfs,cal=euler]{mathalfa}
\usepackage{xfrac}
\usepackage[unicode,hidelinks,bookmarks=false]{hyperref}
\newcommand{\ie}{i.e.\ }
\newcommand{\cf}{cf.\ }
\newcommand{\resp}{respectively\ }
\newcommand{\Subsection}{\subsection}
\providecommand{\nobreakdash}{\nobreak\hbox{-}}
\setlength{\emergencystretch}{2em}
\multlinegap0pt
\usepackage{mathtools}
\usepackage[shortlabels]{enumitem}
\usepackage{amssymb}
\usepackage{braket}
\usepackage{siunitx}
\usepackage{float}
\usepackage{tensor}
\usepackage{csquotes}
\newtheorem{proposition}{Proposition}
\newcommand{\N}{\mathbb{N}}
\newcommand{\R}{\mathbb{R}}
\renewcommand{\dim}{\mathrm{dim}}
\title{\Large\textbf{Dimensional Analysis is a Gauge Theory}}
\author{Benjamin Muntz}
\date{Source: arXiv:2607.21695v1 (CC BY 4.0)}
\begin{document}
\maketitle

\noindent\emph{Source and licence.} \Large\textbf{Dimensional Analysis is a Gauge Theory}. Version arXiv:2607.21695v1 (CC BY 4.0), distributed by arXiv under the Creative Commons Attribution 4.0 International licence, http://creativecommons.org/licenses/by/4.0/, as recorded in the arXiv licence metadata for this exact version and shown on the abstract page https://arxiv.org/abs/2607.21695v1. Full licence notice: \texttt{licenses/arxiv-2607.21695-CC-BY-4.0.txt}.\par\smallskip

\noindent\textit{Editorial scope.} Counted unit: the article's proposition eq:isomorphism (source0.tex lines 301--317). The article's complete original proof of it is delivered with it (source0.tex lines 319--330, one \texttt{proof} environment of 12 lines). No mathematics is written, repaired or re-derived: the statement and the proof are the article's own lines, in the article's own notation, and no retained line is substituted or re-flowed.  No further source-local context is needed: every cross-reference of every retained line resolves inside the two delivered spans.  Nothing was omitted from any retained span, so the package carries no omission record.  No retained line contains a citation, so the wrapper carries no bibliography and no bibliography entry.  No retained line contains a figure environment, an image-inclusion command or drawing code, so no image is delivered and the package declares no asset dependency.  Editorial formatting only, in addition: an AMS \texttt{amsart} preamble that restates the article's own declarations and macros used by the retained lines, namely the environments \texttt{proposition} on separate counters, together with the standard packages those lines need.  The retained environments print in this document as Proposition 1 in order of appearance, whereas the article prints the same items with the numbering of its own full document.  The complete native source of the article as distributed in the arXiv e-print for this exact version is bundled unchanged as \texttt{sources/205886/source0.tex}.
\begin{proposition}\label[proposition]{prop:scalebundleisomorphism}
    Let $P\to M$ be the scale bundle and $w_1,\dots,w_n\in \mathfrak{g}^*$ for some $n\in\N$. The map
    \begin{equation}\label{eq:isomorphism}
    \begin{split}
        \mathsf{\Phi}\colon\quad P & \to \mathcal{L}_{w_1}^+\times_M\cdots \times_M\mathcal{L}_{w_n}^+\\
        p & \mapsto ([p,1],\dots,[p,1]) 
    \end{split}
    \end{equation}
    is a smooth bundle morphism. It is an isomorphism iff $\{w_1,\dots,w_n\}$ forms a complete basis of $\mathfrak{g}^*$. Equivalently, iff the character map
    \begin{equation}
    \begin{split}
        \chi\colon\quad G & \to \R^n_+\\
        g & \mapsto (\rho_{w_1}(g),\dots,\rho_{w_n}(g))
    \end{split}
    \end{equation}
    is a Lie group isomorphism.
\end{proposition}

\begin{proof}
    Compute
    \begin{equation}
    \begin{split}
        \mathsf{\Phi}(pg) &= ([pg,1],\dots,[pg,1])\\
        &= ([p,\rho_{w_1}(g)],\dots,[p,\rho_{w_n}(g)])\\
        &= \mathsf{\Phi}(p)\cdot \chi(g)
    \end{split}
    \end{equation}
    where the last equality indicates the character acting diagonally from the right. So $\mathsf{\Phi}$ is $G$-equivariant after transporting the $G$-action through $\chi$. If $\chi$ is an isomorphism, then $\mathsf{\Phi}$ becomes a bundle morphism between two principal $G$-bundles.\\
    Next we need to show that this is equivalent to saying $\{w_1,\dots,w_n\}$ must be a complete basis of $\mathfrak{g}^*$. Consider therefore the map $W\colon \mathfrak{g}\to \R^n$ with $W(X)\coloneqq (w_1(X),\dots,w_k(X))$ such that $\chi = \exp\circ W\circ \log$. Since $\log$ and $\exp$ are diffeomorphisms, $\chi$ is an isomorphism iff $W$ is a linear isomorphism. But that is only the case iff $\{w_1,\dots,w_n\}$ is a complete basis of $\mathfrak{g}^*$. (And so necessarily $n=k=\dim\,\mathfrak{g}^*$.) This completes the proof.
\end{proof}

\end{document}
